\chapter{MARCO TEÓRICO}

El presente capítulo desarrolla los conceptos fundamentales que sustentan la arquitectura del proyecto y explican por qué cada módulo propuesto aporta a la reducción del retraso promedio por vehículo en intersecciones urbanas. El foco es \textit{arquitectural}: qué módulos componen el sistema, qué principios teóricos utiliza cada uno y qué resultados espera producir; el capítulo siguiente describe el flujo operativo.

En particular, se definen los elementos teóricos que habilitan: (i) la construcción del entorno de simulación en SUMO \citep{SUMO_Netconvert_OSM,OSM_About}, (ii) los controles semafóricos de referencia (tiempo fijo y actuado) contra los que se mide toda mejora \citep{FHWA_Timing_Manual_Chapter3,HCM2000_or_2010}, (iii) el control semafórico adaptativo mediante Aprendizaje por Refuerzo (RL) con algoritmos tabulares y de \textit{Deep RL} \citep{SuttonBarto2018,Mnih2015_DQN,Schulman2017_PPO}, (iv) la instrumentación de métricas de movilidad urbana \citep{FHWA_Toolbox_Definitions,HCM2000_or_2010} y (v) la coordinación distribuida multiagente entre semáforos para optimizar a nivel de corredor \citep{Yu2021_MAPPO}. Los cinco elementos están implementados en esta fase: el RL tabular en el módulo U2, y el \textit{Deep RL} y la coordinación multiagente en el módulo U4, con IPPO; de la comparación algorítmica quedan pendientes SARSA y DQN.

\section{Arquitectura teórica del sistema (visión por módulos)}

\begin{figure}[H]
    \centering
    \includegraphics[width=\textwidth]{Diagrama} % images/Diagrama.png (Overleaf distingue mayusculas)
    \caption{Diagrama de diseño del proyecto: U1 (preprocesamiento), U2 (RL),
    U3 (simulación/medición en SUMO) y U4 (coordinación distribuida).}
    \label{fig:diseno}
\end{figure}

\textbf{U1. Preprocesamiento y generación del entorno SUMO.} Se encarga de producir una red SUMO utilizable a partir de una intersección o pequeño corredor urbano, sea por conversión de cartografía (OSM $\rightarrow$ \texttt{.net.xml}) o por descripción directa de nodos y vías, junto con la definición de flujos y rutas (\texttt{.rou.xml}) y la configuración base de semáforos (\texttt{.add.xml}). Teóricamente, se apoya en la \textit{microsimulación} vehicular y en la representación de grafos viales donde los nodos son intersecciones y las aristas son tramos con uno o más carriles y restricciones de giro \citep{SUMO_Netconvert_OSM,OSM_About}.

\textbf{U2. Control semafórico con Aprendizaje por Refuerzo (RL).} Modela cada semáforo como un agente que observa el estado del entorno (colas por aproximación y fase verde activa), decide si mantener la fase o avanzar a la siguiente y recibe una recompensa que guía el aprendizaje. Se consideran algoritmos tabulares (Q-learning, SARSA) y de \textit{Deep RL} (DQN, PPO). Los elementos teóricos clave son: espacio de estados, espacio de acciones, función de recompensa, política, episodio e interacción agente--entorno \citep{SuttonBarto2018,Mnih2015_DQN,Schulman2017_PPO}. La formulación concreta implementada se detalla en el Capítulo~III.

\textbf{U3. Simulación y medición en SUMO (TraCI).} SUMO provee la dinámica del tráfico a nivel de vehículo; mediante TraCI se leen las variables de estado (vehículos detenidos por aproximación, tiempo de espera acumulado, fase activa) y se aplican las acciones decididas por los agentes. Teóricamente, este módulo sostiene la \textit{retroalimentación} RL: estados y recompensas se computan a partir de la simulación paso a paso. Las métricas de evaluación (retraso, longitud de colas, \textit{throughput}, paradas, emisiones) se calculan sobre las mismas salidas por vehículo para todos los controles comparados, con base en la literatura de ingeniería de transporte \citep{TraCI_Docs,TraCI_Whitepaper,SUMO_Emissions,FHWA_Timing_Manual_Chapter3,FHWA_Toolbox_Definitions,HCM2000_or_2010}.

\textbf{U4. Coordinación distribuida entre semáforos (multiagente).} Extiende el RL local hacia la cooperación entre intersecciones vecinas mediante una capa ligera de mensajería (fase actual, tiempo en la fase, vehículos que descargan hacia el vecino, colas en la aproximación y en el enlace compartido), un \textit{estado extendido} (información propia más un resumen de cada vecino) y una recompensa con término de penalización por \textit{spillback}. En términos teóricos, se apoya en \textit{RL multiagente} y \textit{sistemas distribuidos} (comunicación con latencia y pérdida de mensajes), con IPPO como algoritmo de aprendizaje \citep{Yu2021_MAPPO}. Se implementa en el Capítulo~III sobre los escenarios de varios cruces, en variantes incrementales que permiten atribuir cada mejora a su causa.

\section{Definiciones y conceptos fundamentales}

\subsection*{SUMO (Simulation of Urban MObility)}
Microsimulador de código abierto que modela el movimiento de vehículos a nivel individual sobre una red vial dirigida. Permite definir redes, rutas, tipos de vehículo y control semafórico, y calcular métricas de desempeño (tiempo de viaje, colas, \textit{throughput}, emisiones) \citep{TraCI_Whitepaper,SUMO_Emissions}.

\subsection*{TraCI (Traffic Control Interface)}
Interfaz cliente–servidor que conecta SUMO con programas externos (p.~ej., Python), habilitando control en tiempo de simulación: lectura de observaciones (colas por carril, fase activa, tiempos de espera) y aplicación de acciones (cambios de fase, duraciones) \citep{TraCI_Docs,TraCI_Whitepaper}.

\subsection*{Control de tiempo fijo y control actuado}
Son los dos controles de referencia contra los que se mide cualquier mejora. El \textbf{control de tiempo fijo} repite un programa cíclico de fases con duraciones predefinidas, independientes del tráfico presente; su calidad depende por completo del ajuste del reparto de verde y, en un corredor, del \textit{desfase} entre intersecciones. El \textbf{control actuado} mantiene las mismas fases pero extiende el verde mientras los detectores siguen registrando vehículos, entre un mínimo y un máximo declarados; reacciona a la demanda instantánea sin aprender ni anticipar; SUMO lo implementa de forma nativa declarando una duración mínima y una máxima por fase \citep{SUMO_TLS_Actuated}. Un baseline de tiempo fijo sin ajustar sobrestima cualquier mejora, por lo que en esta tesis el programa fijo se afina por barrido y el actuado se configura con las mismas restricciones de verde que el agente \citep{FHWA_Timing_Manual_Chapter3,HCM2000_or_2010}.

\subsection*{Aprendizaje por Refuerzo (RL)}
Paradigma de aprendizaje donde un agente interactúa con un entorno para maximizar la recompensa acumulada. Elementos: \textit{estado} ($s$), \textit{acción} ($a$), \textit{recompensa} ($r$), \textit{política} ($\pi$) y \textit{episodio}. Un episodio es una corrida completa del entorno, desde su estado inicial hasta un criterio de fin; en esta tesis equivale a una hora de demanda simulada sobre el mismo escenario, más el tiempo necesario para que la red se vacíe. El agente conserva lo aprendido de un episodio al siguiente, de modo que los episodios sucesivos no son observaciones independientes: no pueden tratarse como réplicas estadísticas, y por eso el desempeño se mide aparte, con la política ya fijada. En control semafórico, la recompensa penaliza la congestión observada; su forma concreta (qué componentes y con qué pesos) es una decisión de diseño que se detalla en el Capítulo~III \citep{SuttonBarto2018}.

\subsection*{Política congelada y evaluación}
Durante el entrenamiento el agente explora, es decir, toma acciones al azar con probabilidad $\varepsilon$, y su tabla de valores cambia dentro de cada episodio. Las métricas de esos episodios describen el proceso de aprendizaje, no el desempeño del controlador. Para medir este último se \textit{congela} la política: se fija $\varepsilon = 0$, se deja de actualizar la tabla y se corre el escenario con esa política estática, del mismo modo que se corren los controles de referencia. Solo esos valores son comparables con el tiempo fijo y el actuado.

\subsection*{Q-learning y SARSA (RL tabular)}
Métodos que almacenan el valor de acción $Q(s,a)$ en una tabla indexada por estados discretos. Q-learning es \textit{off-policy}: actualiza sus valores hacia la mejor acción disponible en el estado siguiente, con independencia de la que se haya ejecutado, lo que permite explorar al azar y aprender igualmente la política óptima. SARSA es \textit{on-policy}: actualiza hacia la acción efectivamente tomada, lo que produce políticas más conservadoras. Ambos requieren \textbf{discretizar} el estado, es decir, agrupar magnitudes continuas en un número finito de niveles. Esa discretización es la que hace viable la tabla y, a la vez, su límite: el número de estados crece de forma multiplicativa con las variables observadas y sus niveles, y cuando el espacio es grande el agente visita solo una fracción de él, de modo que la política elige entre valores apenas actualizados. En esta tesis Q-learning es el algoritmo implementado y SARSA queda como comparación pendiente \citep{SuttonBarto2018}.

\subsection*{DQN y PPO (Deep RL)}
Sustituyen la tabla por una red neuronal, lo que permite generalizar entre estados parecidos en lugar de aprender cada uno por separado. DQN aproxima $Q(s,a)$ con una red profunda (\textit{value-based}) y estabiliza el aprendizaje con dos mecanismos: una memoria de experiencias pasadas que se remuestrean y una red objetivo que se actualiza con retraso. PPO optimiza directamente la política (\textit{policy-based}) limitando cuánto puede cambiar de una actualización a la siguiente, mediante el recorte de la razón entre la política nueva y la anterior. Son los candidatos naturales cuando la representación tabular toca su techo \citep{Mnih2015_DQN,Schulman2017_PPO,SB3_Docs}.

\subsection*{Agentes independientes y RL multiagente: IPPO / MAPPO}
Cuando hay varias intersecciones, la configuración más simple es la de \textbf{agentes independientes}: cada semáforo tiene su propia política, observa solo su entorno local y recibe solo su recompensa local, sin intercambiar información con los demás. Cada agente percibe entonces un entorno no estacionario, porque los vecinos también cambian su comportamiento mientras aprenden. En \textbf{IPPO} cada agente entrena su política de forma independiente, con posibles pesos compartidos, y su crítico solo ve la observación propia; \textbf{MAPPO} introduce un crítico centralizado que observa el estado conjunto durante el entrenamiento para dar estabilidad en entornos cooperativos. En esta tesis el módulo U2 usa agentes tabulares independientes en corredor y red, y el módulo U4 usa IPPO con un actor y un crítico compartidos por todos los semáforos, donde la coordinación entra por la observación (los mensajes del vecino) y por la recompensa, no por un crítico centralizado; MAPPO queda como extensión \citep{Yu2021_MAPPO}.

\subsection*{Métricas de movilidad}
\textbf{Retraso (\textit{delay})}: diferencia entre el tiempo de viaje real y el que tomaría el mismo recorrido a la velocidad deseada, sin interferencia de otros vehículos ni del semáforo; en SUMO corresponde al atributo \texttt{timeLoss} de cada viaje e incluye frenadas, detenciones y aceleraciones. Es la métrica primaria de esta tesis. \textbf{Espera detenida}: segundos en que el vehículo permanece prácticamente parado (velocidad $\leq 0.1$~m/s; en SUMO, \texttt{waitingTime}). Es una parte del retraso, no su equivalente: un control puede reducir la espera detenida y aumentar las paradas sin mejorar el retraso. \textbf{Longitud de cola}: número de vehículos detenidos, sumado sobre todas las aproximaciones semaforizadas del escenario; se reporta como promedio por segundo de simulación y como máximo instantáneo. \textbf{\textit{Throughput}}: vehículos que completan su recorrido dentro de la hora simulada; solo discrimina entre controles cuando la demanda supera la capacidad. \textbf{Paradas por vehículo}: número de veces que la velocidad del vehículo cae por debajo del umbral de detención (en SUMO, \texttt{waitingCount}). \textbf{Emisiones}: estimadas por el modelo de SUMO a partir de aceleraciones y velocidades, y acumuladas por vehículo \citep{FHWA_Timing_Manual_Chapter3,FHWA_Toolbox_Definitions,SUMO_Emissions,HCM2000_or_2010}.

Conviene precisar el alcance de la analogía con la literatura de transporte: el \textit{control delay} del HCM y de la FHWA aísla el retraso atribuible al semáforo, mientras que el \texttt{timeLoss} de SUMO incluye además el tiempo perdido en arranques, giros y variaciones de conducción. Por eso un cambio porcentual medido sobre \texttt{timeLoss} subestima el efecto del control: parte del retraso medido no depende de él.

\subsection*{\textit{Spillback} y ondas verdes}
\textit{Spillback}: cola que rebasa la intersección y bloquea enlaces adyacentes; en esta tesis se mide como el tiempo en que la cola detenida de un carril del enlace entre dos cruces ocupa al menos el 90\,\% de sus plazas, y penalizarlo explícitamente en la recompensa es una de las variantes del módulo U4. \textit{Ondas verdes}: alineación temporal de las fases de intersecciones sucesivas, mediante un desfase igual al tiempo de viaje entre ellas, de modo que un pelotón que sale de la primera encuentre verde en las siguientes. Es la forma clásica de coordinar un corredor con control de tiempo fijo, y por eso constituye una referencia exigente para cualquier agente que decida de forma puramente local \citep{FHWA_Timing_Manual_Chapter3,HCM2000_or_2010}.

\subsection*{Reproducibilidad y validez}
La comparación entre controles exige fijar de antemano el protocolo experimental. En esta tesis: las semillas de \textbf{entrenamiento} y de \textbf{evaluación} son disjuntas, de modo que la política no se mide sobre las mismas realizaciones con las que aprendió; la evaluación se hace con la \textbf{política congelada} y con las \textbf{mismas semillas} para todos los controles, lo que permite comparar por pares; la \textbf{unidad experimental} es la corrida de evaluación y no el episodio de entrenamiento, porque los episodios sucesivos comparten la tabla de valores y no son independientes; y se reporta media con \textbf{intervalo de confianza al 95\%}, cambio porcentual con su intervalo, prueba pareada de significancia y tamaño de efecto. Además se documentan las versiones del simulador y de las bibliotecas y se conservan los archivos de resultados de cada corrida.

\section{Relación con objetivos y resultados esperados}

La arquitectura teórica se alinea con los objetivos específicos: (U1) producir una red y rutas SUMO reproducibles; (U2) diseñar y entrenar agentes RL (tabulares y \textit{Deep RL}); (U3) medir con rigor las métricas de movilidad frente a los controles de referencia; (U4) evaluar la coordinación distribuida y su beneficio a nivel de corredor (tiempo extremo a extremo, tasa de \textit{spillback}). El resultado esperado es una reducción del retraso y de las colas frente al control fijo y al actuado, con mejoras adicionales al activar la coordinación. El Capítulo~\ref{cap:resultados} contrasta esa expectativa con lo medido en dos tiempos. Con la representación tabular la reducción solo se obtiene en el escenario más simple, el control actuado resulta mejor que el agente en los cinco casos, y los agentes sin comunicación quedan por detrás de un programa fijo coordinado; esos tres resultados son los que motivaron el módulo U4. Con aproximación de funciones, recompensa de retraso y mensajería entre vecinos, el agente cumple la meta en el cruce, supera al actuado en el corredor e iguala a la onda verde en la red \citep{FHWA_Timing_Manual_Chapter3,HCM2000_or_2010,Yu2021_MAPPO}.
